% ============================================================
%  Semana 1 · Clase 1 — Encuentro Inicial (Refrigerio)
%  Curso: Análisis y Visualización de Datos para la Toma de Decisiones
%
%  Compilar (dos veces) desde la raíz del curso (20262/):
%    TEXINPUTS=".:plantilla:$TEXINPUTS" xelatex clases/semana01-clase01.tex
%    TEXINPUTS=".:plantilla:$TEXINPUTS" xelatex clases/semana01-clase01.tex
% ============================================================
\documentclass[aspectratio=169,11pt]{beamer}

\usepackage{beamerthemeesumer}
\usetheme{esumer}

\usepackage[spanish]{babel}
\usepackage{booktabs}
\usepackage{graphicx}

% ---------------- DATOS DE LA CLASE ----------------
\title{Encuentro Inicial}
\subtitle{Bienvenida al curso}
\author{Docente: [Nombre Apellido]}
\date{Sep 1 al 7}

\setmodulo{Módulo 1: Explorar (Fundamentos y Consulta de Datos)}
\setsemana{1}
\setclase{1}
\setfechaclase{Sep 1 al 7}
\setmodalidad{Híbrido}
% -----------------------------------------------------

\begin{document}

{
\setbeamertemplate{footline}{}
\begin{frame}[plain]
  \titlepage
\end{frame}
}

\begin{frame}{Agenda}
  \begin{enumerate}
    \item Bienvenida
    \item Presentación (docente y estudiantes)
    \item Panorama del curso: los 5 módulos
    \item Cronograma y modalidades
    \item Metodología y evaluación
    \item Logística: herramientas y canales
    \item Cierre + refrigerio
  \end{enumerate}
\end{frame}

\section{Bienvenida}
\begin{frame}{¡Bienvenidos!}
  \begin{block}{Análisis y Visualización de Datos para la Toma de Decisiones}
    \begin{itemize}
      \item Un curso práctico: de los datos crudos a una decisión de negocio sustentada.
      \item 15 semanas · 89 horas · modalidad híbrida (virtual + presencial).
      \item Cierra con un \textbf{proyecto integrador}: un tablero (dashboard) empresarial propio.
    \end{itemize}
  \end{block}
\end{frame}

\section{Presentación}
\begin{frame}{Presentación del docente}
  \begin{itemize}
    \item Nombre, trayectoria y experiencia relacionada con datos/analítica.
    \item Expectativas y estilo de trabajo en el curso.
    \item Canal de contacto para dudas.
  \end{itemize}
\end{frame}

\begin{frame}{Presentación de los estudiantes}
  \begin{alertblock}{Dinámica rápida}
    En parejas o grupos pequeños: nombre, a qué se dedican y una decisión de su
    trabajo/estudio que \emph{les gustaría} poder tomar con mejores datos.
  \end{alertblock}
\end{frame}

\section{Panorama del curso}
\begin{frame}{Los 5 módulos del curso}
  \begin{enumerate}
    \item \textbf{Explorar} — Fundamentos y consulta de datos (Semanas 1–4).
    \item \textbf{Construir} — Limpieza y gestión de datos con Pandas (Semanas 5–8).
    \item \textbf{Experimentar} — Business Intelligence con Power BI (Semanas 9–10).
    \item \textbf{Innovar} — Analítica predictiva y storytelling (Semanas 11–14).
    \item \textbf{Transformar} — Dashboard empresarial aplicado (Semana 15).
  \end{enumerate}
\end{frame}

\begin{frame}{El hilo conductor: un proyecto integrador}
  \begin{itemize}
    \item Desde la Clase 2 se define un problema de negocio y una pregunta analítica.
    \item Cada módulo agrega una capa: estadística $\to$ limpieza $\to$ visualización $\to$ predicción $\to$ storytelling.
    \item El cierre (Semana 15) es la sustentación del tablero final ante un comité.
  \end{itemize}
\end{frame}

\section{Cronograma}
\begin{frame}{Cronograma y modalidades}
  \begin{table}
    \small
    \begin{tabular}{@{}lll@{}}
      \toprule
      \textbf{Semanas} & \textbf{Módulo} & \textbf{Hitos presenciales} \\
      \midrule
      1 -- 4  & Explorar   & MasterClass semana 2 y 4 \\
      5 -- 8  & Construir  & MasterClass semana 7 · Bootcamp semana 8 \\
      9 -- 10 & Experimentar & MasterClass semana 10 \\
      11 -- 14 & Innovar   & MasterClass semana 13 · Bootcamp semana 14 \\
      15      & Transformar & Semana de la Innovación (cierre) \\
      \bottomrule
    \end{tabular}
  \end{table}
  \vspace{0.3cm}
  \small El detalle semana a semana está en el cronograma del curso (\texttt{cronograma.md}).
\end{frame}

\section{Metodología y evaluación}
\begin{frame}{Cómo se evalúa}
  \begin{itemize}
    \item Evidencia/producto en la mayoría de las sesiones (talleres, laboratorios, retos).
    \item Un producto integrador por módulo que se acumula hacia el tablero final.
    \item Sustentación final en la Semana 15 (Semana de la Innovación Multidisciplinar).
  \end{itemize}
  \begin{alertblock}{Pendiente por confirmar}
    Ponderación exacta / rúbrica de evaluación — completar según los criterios institucionales.
  \end{alertblock}
\end{frame}

\section{Logística}
\begin{frame}{Herramientas del curso}
  \begin{itemize}
    \item \textbf{Python:} Google Colab.
    \item \textbf{Bases de datos:} DBeaver + SQLite.
    \item \textbf{Visualización / BI:} Power BI Desktop.
    \item \textbf{Storytelling:} Canva / Google Slides.
    \item \textbf{Opcional:} GitHub para control de versiones del proyecto.
  \end{itemize}
\end{frame}

\begin{frame}{Antes de la próxima clase}
  \begin{itemize}
    \item Crear/activar una cuenta de Google (para Colab).
    \item Instalar DBeaver (se usará desde la Semana 4).
    \item Traer un caso o problema real de su entorno laboral/académico para el proyecto integrador.
  \end{itemize}
\end{frame}

\section{Cierre}
\begin{frame}{Cierre}
  \begin{itemize}
    \item Próxima sesión: \textbf{MasterClass (Nivelación)} — panorama del ecosistema de datos (Semana 2).
    \item Dudas y comentarios.
    \item ¡Vamos por el refrigerio!
  \end{itemize}
\end{frame}

{
\setbeamertemplate{footline}{}
\begin{frame}[plain]
  \begin{tikzpicture}[remember picture,overlay]
    \fill[esumerNavy] (current page.south west) rectangle (current page.north east);
    \node at (current page.center) {\color{white}\Huge\bfseries ¡Gracias!};
  \end{tikzpicture}
\end{frame}
}

\end{document}
